% MA 214 Discussion 2 -- covers Lectures 2 and 3.
%
% Student handout. Page 1 is the concept review; the question set runs from page 2.
% Compile from inside this folder:  pdflatex Discussion2.tex

\documentclass[11pt]{article}
\usepackage[margin=1in]{geometry}
\usepackage{parskip}
\usepackage{fancyhdr}
\usepackage{array}
\usepackage{booktabs}
\usepackage{enumitem}
\usepackage{amsmath}
\usepackage{tabularx}
\usepackage{listings}

\pagestyle{fancy}
\fancyhf{}
\cfoot{\thepage}
\rhead{Discussion 2}
\lhead{MA 214: Applied Statistics}
\renewcommand{\headrulewidth}{.4mm}

\newcommand{\blankline}[1]{\underline{\hspace{#1}}}
\newcolumntype{L}{>{\raggedright\arraybackslash}X}

\lstset{
  basicstyle=\ttfamily\small,
  columns=fullflexible,
  frame=single,
  framerule=0.4pt,
  breaklines=true,
  showstringspaces=false,
  aboveskip=8pt,
  belowskip=8pt
}

% Section header: bold title over a hairline rule.
\newcommand{\parttitle}[1]{%
  \par\medskip\noindent
  {\large\bfseries #1}\par
  \vspace{-.5em}\noindent\rule{\linewidth}{0.4pt}\par\smallskip}

% Writing space.
\newcommand{\answerspace}[1]{\vspace{#1}}

% Flush-left question list: the label runs in at the margin and the body wraps
% back to the margin, so nothing is pushed far to the right.
\newlist{questions}{enumerate}{1}
\setlist[questions]{label=\textbf{Question \arabic*.}, wide=0pt, itemsep=1.4em, topsep=.8em}

\newlist{parts}{enumerate}{1}
\setlist[parts]{label=\textbf{(\Alph*)}, leftmargin=1.6em, labelsep=.45em,
  itemsep=.7em, topsep=.5em}

\begin{document}

\noindent\textbf{Name:}\blankline{2.3in}\hfill\textbf{BUID:}\blankline{1.6in}

\begin{center}
  {\Large\bfseries Discussion 2: One Estimation Principle, Three Models}
\end{center}

\noindent\textbf{Goals:} \textbf{(1)} write a likelihood and log-likelihood for a given model,
\textbf{(2)} find an MLE by differentiating, \textbf{(3)} compute least squares estimates by hand
and interpret them with units, \textbf{(4)} explain why least squares \emph{is} maximum
likelihood, \textbf{(5)} find statistical bugs in R code that runs without error.

\parttitle{Part 1: Concept Review}

{\small

\noindent Everything in Part 2 can be answered from this page. Keep it in front of you.

\begin{enumerate}[label=\textbf{\arabic*.}, ref=\arabic*, wide=0pt, itemsep=.55em, topsep=.5em]

\item \textbf{A statistical model, and the MLE recipe.} ``Trip duration is 14.2 minutes'' is not a
statistical model: it states no randomness and has no parameter to estimate. A statistical model
specifies a \textbf{distribution} for the data, with unknown parameters. To estimate one:
\textbf{(a)} write the \textbf{likelihood} $L(\theta)$, the probability or density of the data you
actually observed, read as a function of $\theta$; \textbf{(b)} take logs,
$\ell(\theta) = \log L(\theta)$, dropping constants free of $\theta$ because they shift $\ell$
without moving where its maximum sits; \textbf{(c)} set $\ell'(\theta) = 0$ and solve for
$\hat\theta$.

\item \textbf{Three models, three MLEs.} Independence across observations is assumed in all three,
which is what lets the likelihood be a product.

\begin{center}
\renewcommand{\arraystretch}{1.15}
\begin{tabular}{llll}
\toprule
\textbf{Data} & \textbf{Model} & \textbf{Likelihood} & \textbf{MLE} \\
\midrule
$y$ of $n$ successes & Binomial$(n,\theta)$ & $\propto \theta^{y}(1-\theta)^{n-y}$ & $\hat\theta = y/n$ \\
$x_1,\dots,x_n > 0$ & Exponential$(\lambda)$ & $\lambda^{n}e^{-\lambda\sum x_i}$ & $\hat\lambda = 1/\bar{x}$ \\
$(x_i, y_i)$ pairs & Normal regression & see item~\ref{cr:mle} & least squares \\
\bottomrule
\end{tabular}
\end{center}

\item\label{cr:model} \textbf{The simple linear regression model.} Written as a statement about a
distribution,
\[
Y_i \mid x_i \;\sim\; \text{Normal}\big(\beta_0 + \beta_1 x_i,\; \sigma^2\big),
\qquad \text{independently across } i .
\]
Three parameters: $\beta_0$, $\beta_1$, $\sigma^2$. Four assumptions ride on that one line: the
mean response is \textbf{linear} in $x$; observations are \textbf{independent}; the
\textbf{variance} $\sigma^2$ is the same at every $x$; the errors are \textbf{normal}.

\item \textbf{Least squares, by hand.} With $S_{xy} = \sum (x_i-\bar{x})(y_i-\bar{y})$ and
$S_{xx} = \sum (x_i-\bar{x})^2$,
\[
\hat\beta_1 = \frac{S_{xy}}{S_{xx}},
\qquad
\hat\beta_0 = \bar{y} - \hat\beta_1\bar{x},
\qquad
\hat{y}_i = \hat\beta_0 + \hat\beta_1 x_i,
\qquad
e_i = y_i - \hat{y}_i ,
\]
and $\text{RSS} = \sum e_i^2$. Note that $\sum e_i = 0$ always, for any fitted model with an
intercept.

\item\label{cr:mle} \textbf{Least squares \emph{is} maximum likelihood.} For the model in
item~\ref{cr:model}, dropping terms free of $\beta_0$ and $\beta_1$,
\[
\ell(\beta_0,\beta_1) \;\propto\; -\frac{1}{2\sigma^2}\sum_{i=1}^{n}\big(y_i - \beta_0 - \beta_1 x_i\big)^2 .
\]
Maximizing a negative multiple of the RSS is the same as minimizing the RSS. The squared error
appears only because the normal density carries $e^{-(\,\cdot\,)^2/2\sigma^2}$.

\item \textbf{Four traps.} Using $\sum e_i$ to judge fit, when it is zero for every model.
Swapping the response and the predictor, since regressing $y$ on $x$ and $x$ on $y$ are different
fits and the second slope is not the reciprocal of the first. Reporting a slope without units,
because ``2.2 more'' is not an interpretation. Interpreting an intercept that sits outside the
data or describes an impossible case.

\end{enumerate}

}

\newpage

\parttitle{Part 2: Question Set}

\noindent\textbf{Scenario.} You are still working with the city bike-share system. Three separate
questions have landed on your desk, and each one needs a different model. All three are answered
the same way: \emph{write down the likelihood, then maximize it.} The third turns out to be the
first two in disguise.

\begin{questions}

%%%%%%%%%%%%%%%%%%%%%%%%%%%%%%%%%%%%%%
\item \textbf{A proportion.} The city audits 20 randomly chosen stations for a broken dock.
\textbf{7 of the 20} have at least one. Let $\theta$ be the probability that a station has a
broken dock.

\begin{parts}
\item Name the distribution and its parameter, and state the assumption that lets you use it
across all 20 stations.
\answerspace{3.5em}

\item Write $L(\theta)$, then $\ell(\theta)$. You may drop the binomial coefficient; say why you
are allowed to.
\answerspace{3.5em}

\item Differentiate, set $\ell'(\theta) = 0$, and solve.
\answerspace{3.5em}

\item Report $\hat\theta$ as a number, and say in one sentence what it estimates, in context.
\answerspace{3.5em}

\item Suppose instead 70 of 200 stations were broken. What is $\hat\theta$ now? Since the estimate
did not change, what \emph{did} change about the likelihood function?
\answerspace{3.5em}
\end{parts}

%%%%%%%%%%%%%%%%%%%%%%%%%%%%%%%%%%%%%%
\item \textbf{A rate.} A rider waits at an empty station until a bike is returned. Five such
waits, in minutes:
\[
2.5, \qquad 1.0, \qquad 4.0, \qquad 0.5, \qquad 2.0
\]
Model each wait as exponential with rate $\lambda$, so $f(x) = \lambda e^{-\lambda x}$ for
$x > 0$, and treat the waits as independent.

\begin{parts}
\item Write $L(\lambda)$ for all five observations as a single expression.
\answerspace{3.5em}

\item Write $\ell(\lambda)$, differentiate, and solve. Express $\hat\lambda$ in terms of $\bar{x}$
before plugging in numbers.
\answerspace{3.5em}

\item Give $\hat\lambda$ as a number, with units.
\answerspace{3em}

\item The mean of an exponential distribution is $1/\lambda$. What is the estimated mean wait, and
why should that not surprise you given (B)?
\answerspace{3.5em}
\end{parts}

%%%%%%%%%%%%%%%%%%%%%%%%%%%%%%%%%%%%%%
\item \textbf{A regression slope.} Five stations. Here $x$ is the number of docks in tens, and $y$
is daily trips started, in tens.

\begin{center}
\renewcommand{\arraystretch}{1.15}
\begin{tabular}{lccccc}
\toprule
\textbf{Station} & A & B & C & D & E \\
\midrule
$x$ (tens of docks) & 6 & 7 & 8 & 9 & 10 \\
$y$ (tens of trips) & 12 & 14 & 15 & 18 & 21 \\
\bottomrule
\end{tabular}
\end{center}

\begin{parts}
\item Compute $\bar{x}$ and $\bar{y}$, then fill in the table and total the last two columns.

\begin{center}
\renewcommand{\arraystretch}{1.6}
\begin{tabular}{cccc}
\toprule
$x_i - \bar{x}$ & $y_i - \bar{y}$ & $(x_i-\bar{x})(y_i-\bar{y})$ & $(x_i-\bar{x})^2$ \\
\midrule
\blankline{.8in} & \blankline{.8in} & \blankline{1.1in} & \blankline{1.1in} \\
\blankline{.8in} & \blankline{.8in} & \blankline{1.1in} & \blankline{1.1in} \\
\blankline{.8in} & \blankline{.8in} & \blankline{1.1in} & \blankline{1.1in} \\
\blankline{.8in} & \blankline{.8in} & \blankline{1.1in} & \blankline{1.1in} \\
\blankline{.8in} & \blankline{.8in} & \blankline{1.1in} & \blankline{1.1in} \\
\midrule
 & & $S_{xy} = $\blankline{.7in} & $S_{xx} = $\blankline{.7in} \\
\bottomrule
\end{tabular}
\end{center}

\item Compute $\hat\beta_1$ and $\hat\beta_0$, and write the fitted equation.
\answerspace{3.5em}

\item Fill in the fitted values and residuals.

\begin{center}
\renewcommand{\arraystretch}{1.55}
\begin{tabular}{cccc}
\toprule
\textbf{Station} & $y_i$ & $\hat{y}_i$ & $e_i = y_i - \hat{y}_i$ \\
\midrule
A & 12 & \blankline{1.1in} & \blankline{1.1in} \\
B & 14 & \blankline{1.1in} & \blankline{1.1in} \\
C & 15 & \blankline{1.1in} & \blankline{1.1in} \\
D & 18 & \blankline{1.1in} & \blankline{1.1in} \\
E & 21 & \blankline{1.1in} & \blankline{1.1in} \\
\bottomrule
\end{tabular}
\end{center}

\item Add up your residuals: what do you get, and is that a coincidence? Then compute the RSS.
\answerspace{4em}

\item Interpret $\hat\beta_1$ in context, \emph{with units}. Then interpret $\hat\beta_0$ in
context and say whether that interpretation is worth reporting.
\answerspace{5em}
\end{parts}

%%%%%%%%%%%%%%%%%%%%%%%%%%%%%%%%%%%%%%
\item \textbf{The bridge.} You have now maximized two likelihoods. In Question 3 you never wrote
one down.

\begin{parts}
\item Using item~\ref{cr:mle} of the concept review, explain in one sentence why \emph{maximizing}
\[
-\frac{1}{2\sigma^2}\sum_{i=1}^{n}\big(y_i - \beta_0 - \beta_1 x_i\big)^2
\]
is the same as \emph{minimizing} the RSS.
\answerspace{3.5em}

\item Which model assumption puts the \emph{square} in least squares? Circle one and justify.
\begin{enumerate}[label=\arabic*., leftmargin=1.8em, itemsep=.15em, topsep=.3em]
\item the errors are independent
\item the errors are normally distributed
\item the errors have constant variance
\item the relationship is linear
\end{enumerate}
\answerspace{3.5em}

\item If the errors were strongly right-skewed instead of normal, would least squares still be the
maximum likelihood estimator? Would $S_{xy}/S_{xx}$ still be \emph{computable}?
\answerspace{3.5em}
\end{parts}

%%%%%%%%%%%%%%%%%%%%%%%%%%%%%%%%%%%%%%
\item \textbf{Find the bugs.} A classmate wants the slope and the RSS for the station data. The
code runs with no error message and prints plausible-looking numbers. It is wrong in
\textbf{three} separate ways.

\begin{lstlisting}[language=R]
fit <- lm(docks ~ trips, data = stations)

b1 <- coef(fit)[1]

rss <- sum(resid(fit))

c(slope = b1, rss = rss)
\end{lstlisting}

\begin{parts}
\item Bug 1 is in the formula. What did it fit instead, and why is the answer not simply the
reciprocal of the slope they wanted?
\answerspace{3.5em}

\item Bug 2 is in the indexing. What does \texttt{coef(fit)[1]} return?
\answerspace{3em}

\item Bug 3 is in the \texttt{rss} line. What number does it print for \emph{any} fitted model
with an intercept, and what should the line be?
\answerspace{3.5em}

\item Rewrite all three lines correctly.
\answerspace{3.5em}

\item The same classmate writes the two sentences below. Mark each \textbf{defensible} or
\textbf{not}, and rewrite the ones that are not.
\begin{enumerate}[label=\textbf{S\arabic*.}, leftmargin=2.2em, itemsep=.25em, topsep=.35em]
\item ``Our residuals sum to zero, so the line fits the stations well.''
\item ``Least squares works here because the errors are normally distributed.''
\end{enumerate}
\answerspace{5em}
\end{parts}

\end{questions}

\end{document}
